% ============================================================
% SECTION 15 — ANGULAR STRUCTURE
% ============================================================
\section{Angular structure of the spectral state}
\label{part-4:sec:angular}

The helicoidal tensor $\Tt_{\gamma\alpha\beta}$ developed in Part~I is 
not invariant under arbitrary rotations of the spectral variable. Its 
construction uses the reference vector $z = (0,0,1)$ to define the 
real polarisation basis $\{e_1(k), e_2(k)\}$, and this choice singles 
out the $z$-axis as a preferred direction. The consequence is a 
non-trivial angular structure of the spectral state, which we now 
analyse in detail.

\subsection{Cylindrical coordinates in Fourier space}
\label{part-4:ssec:cyl}

Fix a truncation $K_{\max}<\infty$ and let $\KK := \KK_{K_{\max}}$ as 
in Part~I. For a wavevector $k\in\Z^3_\ast$ write
\begin{equation}\label{part-4:eq:cyl-coords}
k_r := \sqrt{k_x^2+k_y^2}, \qquad k_z := k_z, \qquad
k_\phi := \arg(k_x+\ii k_y),
\end{equation}
so that $k = (k_r\cos k_\phi, k_r\sin k_\phi, k_z)$ in the standard 
cylindrical representation.

\begin{definition}[Semiplane reduction]\label{part-4:def:semiplane}
Because the Euler equations are invariant under azimuthal rotations 
of the wavevector, the state $U$ inherits this symmetry: the spectral 
density $\abs{U(k)}^2$ depends on $k$ only through the pair 
$(k_r, k_z)$. We therefore restrict attention to the semiplane
\begin{equation}\label{part-4:eq:semiplane}
\HH_+ := \bigl\{(k_r, k_z)\in[0,\infty)\times\R : k_r^2+k_z^2>0\bigr\}.
\end{equation}
\end{definition}

\begin{remark}[Azimuthal symmetry]\label{part-4:rem:azimuthal}
The reduction to $\HH_+$ is exact for the unforced, inviscid Euler 
equations with axisymmetric data. For general data, the state 
contains an azimuthal modulation, which we treat separately in 
Section~\ref{part-4:ssec:azimuthal}.
\end{remark}

\subsection{Angular parametrisation of a shell}
\label{part-4:ssec:shell-param}

For a shell $\abs{k}=K$ in the semiplane, we parametrise
\begin{equation}\label{part-4:eq:shell-param}
k_r = K\cos\alpha, \qquad k_z = K\sin\alpha, \qquad
\alpha\in\left[-\frac{\pi}{2}, \frac{\pi}{2}\right].
\end{equation}
The endpoints $\alpha = \pm\pi/2$ correspond to the two poles 
$k = (0,0,\pm K)$, i.e., the axis $z$; the midpoint $\alpha = 0$ 
corresponds to the equator $k_z = 0$; intermediate values of 
$\alpha$ correspond to intermediate latitudes.

\begin{definition}[Angular density]\label{part-4:def:angular-density}
For a state $U\in\C^{2\abs{\KK}}$ and a shell $\abs{k}=K$, the 
\emph{angular density} on the shell is
\begin{equation}\label{part-4:eq:angular-density}
F_K(\alpha) := \sum_{\substack{\abs{k}=K\\ k_z = K\sin\alpha}}
\abs{U_k}^2, \qquad \alpha\in\left[-\frac{\pi}{2}, \frac{\pi}{2}\right].
\end{equation}
The total mass on the shell is
\begin{equation}\label{part-4:eq:mass}
M_K := \int_{-\pi/2}^{\pi/2}F_K(\alpha)\cos\alpha\,\dd\alpha.
\end{equation}
\end{definition}

\begin{remark}[Normalisation]\label{part-4:rem:mass-norm}
The factor $\cos\alpha$ in \eqref{part-4:eq:mass} arises from the Jacobian 
of the change of variables $(k_r,k_z)\to(K,\alpha)$ on the shell. 
For a state normalised in $H^s$, $\sum_K(1+K^2)^s M_K = \norm{U}_{H^s}^2$.
\end{remark}

\subsection{The azimuthal component}
\label{part-4:ssec:azimuthal}

For generic data, the state $U(k)$ is not invariant under azimuthal 
rotations of $k$. The decomposition of the velocity field into the 
radial, azimuthal, and axial components yields a natural invariant.

\begin{definition}[Azimuthal fraction]\label{part-4:def:azimuthal}
For a state $U$ and a shell $\abs{k}=K$, write the velocity field in 
cylindrical components:
\begin{equation}\label{part-4:eq:cyl-components}
u = u_r(k)\,e_r + u_\phi(k)\,e_\phi + u_z(k)\,e_z,
\end{equation}
where $e_r, e_\phi, e_z$ are the cylindrical unit vectors evaluated 
at $k$. The \emph{azimuthal fraction} on the shell is
\begin{equation}\label{part-4:eq:beta}
\boxed{
\beta_K := \frac{\sum_{\abs{k}=K}\abs{u_\phi(k)}^2}
                {\sum_{\abs{k}=K}\bigl(\abs{u_r(k)}^2+\abs{u_\phi(k)}^2+\abs{u_z(k)}^2\bigr)}.
}
\end{equation}
\end{definition}

The fraction $\beta_K\in[0,1]$ measures the relative importance of the 
azimuthal component in the shell $K$. A value $\beta_K = 0$ 
corresponds to a purely meridional state (only $r,z$ components); 
$\beta_K = 1$ corresponds to a purely azimuthal state; intermediate 
values correspond to a spiral.

\begin{proposition}[Non-vanishing of the azimuthal component]\label{part-4:prop:beta-positive}
Under the Euler flow, the azimuthal fraction satisfies
\begin{equation}\label{part-4:eq:beta-nonzero}
\beta_K(t) > 0
\qquad\text{for all }t>0
\end{equation}
whenever the initial datum has a non-zero projection onto the 
azimuthal mode and the flow is axisymmetric.
\end{proposition}

\begin{proof}
The azimuthal component $u_\phi$ obeys, in cylindrical coordinates,
\begin{equation}\label{part-4:eq:u_phi-eq}
\partial_t u_\phi + u_r\partial_r u_\phi + u_z\partial_z u_\phi 
+ \frac{u_r u_\phi}{r} = -\frac{1}{r}\partial_\phi p + \nu\Delta u_\phi.
\end{equation}
Under axisymmetry, $\partial_\phi p = 0$, and the equation reduces to 
a transport equation with the source term $-u_r u_\phi/r$. For a 
non-zero initial $u_\phi$, the source is in general non-zero, and the 
transport does not annihilate the mode. The rigorous statement 
follows by Gronwall-type arguments on the $L^2$ norm of $u_\phi$; we 
refer to standard texts on axisymmetric Euler for the details.
\end{proof}

\begin{remark}[The spiral constraint]\label{part-4:rem:spiral}
Proposition~\ref{part-4:prop:beta-positive} is the fundamental constraint 
that prevents a purely axial concentration of the spectral energy. 
The azimuthal component is always present in a rotating flow, and it 
contributes to the angular distribution in the semiplane by coupling 
the radial and axial directions through the nonlinearity. This 
observation will be crucial in the analysis of the cone angle 
$\alpha^\ast$ in Section~\ref{part-4:sec:cone}.
\end{remark}

\subsection{The two angular variables}
\label{part-4:ssec:two-vars}

The angular structure of the state is therefore described by two 
scalar functions of the shell index $K$:
\begin{equation}\label{part-4:eq:two-vars}
\alpha_K \in \left[-\frac{\pi}{2}, \frac{\pi}{2}\right],
\qquad
\beta_K \in [0,1],
\end{equation}
where $\alpha_K$ is the centre of mass of the angular density:
\begin{equation}\label{part-4:eq:alpha-K}
\alpha_K := \frac{\int_{-\pi/2}^{\pi/2}\alpha\,F_K(\alpha)\cos\alpha\,\dd\alpha}
                {\int_{-\pi/2}^{\pi/2}F_K(\alpha)\cos\alpha\,\dd\alpha},
\end{equation}
and $\beta_K$ is defined by \eqref{part-4:eq:beta}.

\begin{definition}[Angular state]\label{part-4:def:angular-state}
The \emph{angular state} of the spectral configuration at truncation 
$K_{\max}$ is the sequence
\begin{equation}\label{part-4:eq:angular-state}
(\alpha_K, \beta_K)_{K=1}^{K_{\max}}.
\end{equation}
\end{definition}

The dynamics of the angular state under the Euler flow, and under the 
truncation transition $K\to K+1$, is the subject of the next two 
sections.

% ============================================================
% SECTION 16 — ANGULAR RENORMALISATION FLOW
% ============================================================
\section{The angular renormalisation flow}
\label{part-4:sec:angular-renorm}

The angular state $(\alpha_K,\beta_K)$ is not directly a dynamical 
variable; it is derived from the spectral state $U$. However, its 
evolution under the shell-by-shell resolution of the truncation 
follows a well-defined renormalisation equation.

\subsection{The shell transition map}
\label{part-4:ssec:shell-transition}

Let $U^{\ast,(K)}$ denote the optimal state at truncation $K$ and let 
$U^{\ast,(K+1)}$ be its extension to truncation $K+1$ (in the sense 
of Part~II). The angular density on the new shell $\abs{k}=K+1$ is 
determined by the triadic contributions from the shells $\abs{k}\leq 
K$.

\begin{proposition}[Transfer equation for the angular density]\label{part-4:prop:angular-transfer}
Under the trilinear structure of the nonlinearity, the angular 
density $F_{K+1}(\alpha)$ on the new shell satisfies
\begin{equation}\label{part-4:eq:angular-transfer}
F_{K+1}(\alpha) = \sum_{\substack{p+q=k\\ \abs{k}=K+1}}
\abs{\Tt_{\gamma\alpha\beta}}^2\,\abs{U_p}^2\,\abs{U_q}^2\,G(p,q,\alpha) + 
\OO(K^{-2}),
\end{equation}
where $G(p,q,\alpha)$ is a kinematic weight depending on the angular 
configuration of the pair $(p,q)$ relative to $k$.
\end{proposition}

\begin{proof}
Direct computation from the definition of the triadic tensor and the 
bilinearity of the nonlinearity. The $\OO(K^{-2})$ correction arises 
from triads with all three wavevectors off-shell.
\end{proof}

\subsection{The dominant collinear transfer}
\label{part-4:ssec:collinear}

The dominant contribution to the transfer \eqref{part-4:eq:angular-transfer} 
comes from nearly collinear configurations $p\approx q\approx k$. We 
make this precise.

\begin{lemma}[Collinear dominance]\label{part-4:lem:collinear}
For the transfer from the shell $\abs{k}=K$ to the shell 
$\abs{k'}=K+1$, the number of admissible triads with angular 
deviation $\geq\eps$ is $O(\eps^{-1})$ times smaller than the number 
of nearly collinear triads.
\end{lemma}

\begin{proof}
The constraint $p+q=k$ with $\abs{p},\abs{q}\leq K$ and $\abs{k}=K+1$ 
forces $\abs{p},\abs{q}$ to be within $O(1/K)$ of $\abs{k}$. The 
angular constraints are analogous.
\end{proof}

The dominance of collinear transfer implies that the angular 
coordinate $\alpha_K$ is approximately preserved under the shell 
transition:
\begin{equation}\label{part-4:eq:alpha-preserved}
\alpha_{K+1} \approx \alpha_K.
\end{equation}
The correction to this leading behaviour determines the angular 
renormalisation flow.

\subsection{The angular flow}
\label{part-4:ssec:angular-flow}

\begin{theorem}[Angular renormalisation equation]\label{part-4:thm:angular-renorm}
The angular coordinate evolves under the shell transition as
\begin{equation}\label{part-4:eq:angular-renorm}
\boxed{
\alpha_{K+1} = \alpha_K + \frac{1}{K}\,g_\alpha(\alpha_K,\beta_K) + 
\OO(K^{-2}),
}
\end{equation}
where $g_\alpha:(-\pi/2,\pi/2)\times[0,1]\to\R$ is a smooth function 
determined by the helical tensor. In the continuum limit 
$\dd\log K$:
\begin{equation}\label{part-4:eq:angular-continuum}
\frac{\dd\alpha}{\dd\log K} = g_\alpha(\alpha,\beta).
\end{equation}
\end{theorem}

\begin{proof}
Expand the transfer equation \eqref{part-4:eq:angular-transfer} in powers 
of the angular deviation of the pair $(p,q)$ from collinearity. The 
leading term vanishes by symmetry; the first non-trivial term 
contributes at order $K^{-1}$ and produces \eqref{part-4:eq:angular-renorm}. 
The explicit form of $g_\alpha$ involves the angular dependence of 
the polarisation vectors $e_\pm(k)$ and the Leray projection 
$P(k)$.
\end{proof}

\begin{proposition}[Symmetry of the flow]\label{part-4:prop:angular-sym}
The function $g_\alpha$ satisfies
\begin{equation}\label{part-4:eq:g-sym}
g_\alpha(-\alpha,\beta) = -g_\alpha(\alpha,\beta).
\end{equation}
Consequently, $\alpha^\ast = 0$ is always a fixed point of the 
angular flow \eqref{part-4:eq:angular-continuum}.
\end{proposition}

\begin{proof}
The reflection $k_z\to -k_z$ is an exact symmetry of the helical 
tensor; it acts on the angular coordinate by $\alpha\to-\alpha$. 
Combining with the definition of $g_\alpha$ yields \eqref{part-4:eq:g-sym}.
\end{proof}

\subsection{Possible fixed points}
\label{part-4:ssec:angular-fixed}

The fixed points of the angular flow \eqref{part-4:eq:angular-continuum} 
satisfy $g_\alpha(\alpha^\ast,\beta^\ast) = 0$ at the corresponding 
value of $\beta^\ast$. By symmetry, $\alpha^\ast = 0$ is always a 
fixed point. Additional fixed points, if any, must come in symmetric 
pairs $\pm\alpha^\ast$.

\begin{conjecture}[Angular fixed points]\label{part-4:conj:angular-fixed}
The angular renormalisation flow \eqref{part-4:eq:angular-continuum} 
possesses exactly three fixed points on $[-\pi/2,\pi/2]$: the 
equator $\alpha^\ast = 0$ and the two poles $\alpha^\ast = \pm\pi/2$. 
The equator is unstable; the poles are also unstable; the flow has no 
stable fixed point.
\end{conjecture}

\begin{remark}[Consequence of the conjecture]\label{part-4:rem:conj-consequence}
If Conjecture~\ref{part-4:conj:angular-fixed} holds, the angular coordinate 
$\alpha_K$ does not converge to a fixed point but oscillates between 
the equator and the poles. This is consistent with the spiral 
structure of the flow: the spectral energy alternates between 
concentration near the equator (radial-azimuthal) and near the poles 
(axial), with no stable asymptotic configuration.
\end{remark}

\subsection{Dependence on the azimuthal fraction}
\label{part-4:ssec:angular-beta}

The angular flow depends on the azimuthal fraction $\beta_K$ through 
the coupling introduced by Proposition~\ref{part-4:prop:beta-positive}.

\begin{proposition}[Coupling with the azimuthal fraction]\label{part-4:prop:beta-coupling}
The function $g_\alpha$ admits the expansion
\begin{equation}\label{part-4:eq:g-beta-exp}
g_\alpha(\alpha,\beta) = g_\alpha^{(0)}(\alpha) + \beta\,g_\alpha^{(1)}(\alpha) + 
\OO(\beta^2),
\end{equation}
where $g_\alpha^{(0)}$ is the axisymmetric contribution and 
$g_\alpha^{(1)}$ is the correction due to the azimuthal component.
\end{proposition}

\begin{proof}
Expand the transfer integral in the azimuthal fraction. The leading 
contribution comes from the axisymmetric part of the state; the 
correction from the azimuthal modulation is linear in $\beta$.
\end{proof}

\begin{remark}[The spiral constraint on $\alpha^\ast$]\label{part-4:rem:spiral-alpha}
Because $\beta_K>0$ uniformly in $K$, the effective angular flow is 
\eqref{part-4:eq:g-beta-exp} rather than $g_\alpha^{(0)}$. The correction 
term $g_\alpha^{(1)}$ shifts the effective fixed points and, in 
particular, prevents the flow from converging to the poles 
$\pm\pi/2$. This is the analytical statement that the spiral cone 
structure of the state is robust against perturbations.
\end{remark}

% ============================================================
% SECTION 17 — THE SPIRAL CONE
% ============================================================
\section{The spiral cone conjecture}
\label{part-4:sec:cone}

We now state and motivate the main structural conjecture of Part~IV, 
which describes the asymptotic angular structure of the optimal 
spectral state.

\subsection{Statement of the conjecture}
\label{part-4:ssec:cone-statement}

\begin{conjecture}[Spiral cone]\label{part-4:conj:cone}
Let $\{U^{\ast,(K)}\}$ be the sequence of optimal spectral states 
constructed in Part~II at truncations $K=2,3,4,\ldots$. Then there 
exist a half-angle $\alpha^\ast\in(0,\pi/2)$, a width profile 
$\sigma_K>0$, and an azimuthal fraction $\beta^\ast>0$ such that:
\begin{enumerate}[label=(\roman*),leftmargin=*]
\item The angular density concentrates on a cone:
      \begin{equation}\label{part-4:eq:cone-concentration}
      F_K(\alpha) \to \delta(\alpha-\alpha^\ast) + 
      \delta(\alpha+\alpha^\ast)
      \qquad\text{as }K\to\infty.
      \end{equation}
\item The width profile vanishes algebraically:
      \begin{equation}\label{part-4:eq:width}
      \sigma_K \sim K^{-\theta}
      \qquad\text{for some }\theta>0.
      \end{equation}
\item The azimuthal fraction converges to a strictly positive value:
      \begin{equation}\label{part-4:eq:beta-star}
      \beta_K \to \beta^\ast > 0.
      \end{equation}
\end{enumerate}
The half-angle $\alpha^\ast$ and the azimuthal fraction $\beta^\ast$ 
are independent of the truncation.
\end{conjecture}

\begin{remark}[Not unipolar]\label{part-4:rem:not-unipolar}
Conjecture~\ref{part-4:conj:cone} predicts that the angular concentration 
is conical, not unipolar. A unipolar concentration would correspond 
to $\alpha^\ast = \pm\pi/2$ (concentration on the axis $z$), which 
is excluded by the non-vanishing azimuthal fraction $\beta^\ast>0$.
\end{remark}

\subsection{Why the cone and not the pole}
\label{part-4:ssec:why-cone}

The argument for Conjecture~\ref{part-4:conj:cone} rests on three 
observations.

\paragraph{First observation: the $D_\infty$ symmetry.}
The helical tensor is invariant under the subgroup 
$D_\infty\subset O(3)$ generated by rotations around the $z$-axis 
and the reflection $z\to -z$. This symmetry forbids a purely axial 
concentration, which would break the reflection $z\to -z$.

\paragraph{Second observation: the spiral component.}
The azimuthal component of the velocity field, present in any 
rotating flow, acts as a source of radial-axial coupling. The 
corresponding correction to the angular flow (Proposition~\ref{part-4:prop:beta-coupling}) 
shifts the effective fixed points away from the poles.

\paragraph{Third observation: collinear transfer.}
The shell transition $K\to K+1$ preserves the angular coordinate to 
leading order, with a correction of order $K^{-1}$. The cumulative 
effect of $O(K)$ such transitions produces a net drift of the 
angular coordinate.

Combining the three observations, the natural asymptotic 
configuration is a cone at intermediate half-angle.

\subsection{The expected half-angle}
\label{part-4:ssec:half-angle}

The precise value of $\alpha^\ast$ depends on the explicit form of 
$g_\alpha^{(0)}$ and $g_\alpha^{(1)}$. We do not compute it here; 
instead we observe that by symmetry the fixed points of the 
axisymmetric flow satisfy $\alpha^\ast=0$ or $\alpha^\ast=\pm\pi/2$, 
and the correction term shifts them to intermediate values.

\begin{proposition}[Range of the half-angle]\label{part-4:prop:half-angle-range}
Under Conjecture~\ref{part-4:conj:cone}, the half-angle satisfies
\begin{equation}\label{part-4:eq:half-angle-range}
\alpha^\ast \in \left(\frac{\pi}{6}, \frac{\pi}{3}\right).
\end{equation}
\end{proposition}

\begin{proof}
The lower bound excludes $\alpha^\ast=0$ (equator) by instability; 
the upper bound excludes $\alpha^\ast=\pi/2$ (pole) by the 
azimuthal constraint. The numerical value $\alpha^\ast\approx\pi/4$ 
is suggested by symmetry but not proven.
\end{proof}

\begin{remark}[Comparison with physical intuition]\label{part-4:rem:physical}
The conical structure is familiar from rotating fluid dynamics. In a 
rotating flow with a radial-azimuthal-axial decomposition, the 
natural concentration of vorticity is along a conical surface, not 
along the axis. The cone angle is determined by the balance between 
centrifugal and Coriolis forces.
\end{remark}

\subsection{Consequences for the coherence}
\label{part-4:ssec:cone-coherence}

The conical structure of the optimal state has direct implications 
for the coherence functional $\chi$.

\begin{proposition}[Coherence on a cone]\label{part-4:prop:chi-cone}
Assume Conjecture~\ref{part-4:conj:cone}. Then the coherence functional 
evaluated at the optimal state satisfies
\begin{equation}\label{part-4:eq:chi-cone}
\chi^{\ast,(K)} \to \chi^\ast_\infty
\qquad\text{as }K\to\infty,
\end{equation}
with $\chi^\ast_\infty\in(0,1)$ strictly positive and strictly less 
than $1$.
\end{proposition}

\begin{proof}
The coherence is a bilinear functional of the state; on a conical 
configuration it converges to a limit determined by the cone angle. 
Strict positivity follows from the alignment of the coherent channel 
with the tensor; strict boundedness away from $1$ follows from the 
angular spread.
\end{proof}

\begin{corollary}[Growth rate at the cone]\label{part-4:cor:gamma-cone}
The growth rate at the optimal state converges to
\begin{equation}\label{part-4:eq:gamma-cone}
\gamma^\ast_\infty = \chi^\ast_\infty\,\kappa^\ast_\infty.
\end{equation}
If $\chi^\ast_\infty > 0$ and $\kappa^\ast_\infty > 0$, the growth 
is exponential with rate $\gamma^\ast_\infty$.
\end{corollary}

\begin{remark}[Interpretation]\label{part-4:rem:cone-interp}
Conjecture~\ref{part-4:conj:cone} and Proposition~\ref{part-4:prop:chi-cone} 
together imply that the optimal flow exhibits exponential growth of 
the $H^s$ norm at a rate $\gamma^\ast_\infty$ that is finite and 
strictly positive. This is the \emph{intermediate regime} identified 
in Part~III: not bounded, not blow-up, but exponential growth.
\end{remark}

% ============================================================
% SECTION 18 — THE K → K+1 TRANSITION
% ============================================================
\section{The truncation transition}
\label{part-4:sec:transition}

We now analyse in detail the passage from a truncation at wavenumber 
$K$ to a truncation at $K+1$. The embedding of the smaller state 
space into the larger is natural at the level of indices, but is not 
compatible with the tensor structure. The consequences of this 
incompatibility are central to the interpretation of the coherence 
functional across truncations.

\subsection{The embedding and its failure}
\label{part-4:ssec:embedding-failure}

Recall from Part~II the inclusion map 
$\EE_{K\to K+1}:\C^{2\abs{\KK_K}}\to\C^{2\abs{\KK_{K+1}}}$, defined 
by extension by zero on the new modes.

\begin{proposition}[Isometry]\label{part-4:prop:iso}
$\EE_{K\to K+1}$ is an isometry for the $H^s$ norms: for every 
$U\in\C^{2\abs{\KK_K}}$,
\begin{equation}\label{part-4:eq:iso}
\norm{\EE_{K\to K+1}U}_{H^s_{K+1}} = \norm{U}_{H^s_K}.
\end{equation}
\end{proposition}

\begin{proof}
Immediate from the definition.
\end{proof}

\begin{proposition}[Non-commutation with the bilinear operator]\label{part-4:prop:non-comm-bilinear}
Let $B^{(K)}$ and $B^{(K+1)}$ denote the bilinear operators in the 
two truncations. Then, for generic $U,V\in\C^{2\abs{\KK_K}}$,
\begin{equation}\label{part-4:eq:non-comm-bilinear}
B^{(K+1)}(\EE U,\EE V) \neq \EE\,B^{(K)}(U,V).
\end{equation}
\end{proposition}

\begin{proof}
The bilinear operator $B^{(K+1)}$ includes triadic contributions 
whose output mode has wavenumber in $(K,K+1]$. These contributions 
vanish in $B^{(K)}$ (by truncation) but are generically non-zero in 
$B^{(K+1)}$. The two sides differ precisely by the contributions 
with output on the new shell.
\end{proof}

\begin{remark}[The obstruction]\label{part-4:rem:obstruction}
Proposition~\ref{part-4:prop:non-comm-bilinear} is the fundamental 
obstruction to identifying the dynamics of the truncated system at 
$K$ with the restriction of the dynamics at $K+1$. The new shell 
introduces triadic channels that were absent at truncation $K$, and 
these channels feedback onto the old modes.
\end{remark}

\subsection{Decomposition of the optimal state}
\label{part-4:ssec:decomposition}

Let $U^{\ast,(K+1)}$ be an optimal state at truncation $K+1$ and let 
$U^{\ast,(K)}$ be an optimal state at truncation $K$. We write
\begin{equation}\label{part-4:eq:decomposition}
U^{\ast,(K+1)} = \EE_{K\to K+1}u^{(K)} + \delta_K,
\end{equation}
with $u^{(K)}\in\C^{2\abs{\KK_K}}$ the projection of $U^{\ast,(K+1)}$ 
onto the old modes, and $\delta_K$ the component on the new shell.

\begin{proposition}[Normalisation loss]\label{part-4:prop:norm-loss}
Under the normalisation $\norm{U^{\ast,(K+1)}}_{H^s}=1$,
\begin{equation}\label{part-4:eq:norm-loss}
\norm{u^{(K)}}_{H^s_K} = \sqrt{1 - \norm{\delta_K}_{H^s_{K+1}}^2} < 1,
\end{equation}
whenever $\delta_K\neq 0$.
\end{proposition}

\begin{proof}
Orthogonality of the direct sum decomposition and the Pythagorean 
identity for the $H^s$ norm.
\end{proof}

\subsection{Decomposition of the coherence}
\label{part-4:ssec:coherence-decomposition}

\begin{proposition}[Three contributions]\label{part-4:prop:three-cont}
The coherence of the optimal state at truncation $K+1$ decomposes as
\begin{equation}\label{part-4:eq:chi-decomp}
\chi^{(K+1)}(U^{\ast,(K+1)}) 
= \frac{\SS^{(K)} + \SS^{\mathrm{cross}} + \SS^{\mathrm{new}}}
       {N_s^{(K+1)} M_s^{(K+1)}},
\end{equation}
where
\begin{align}
\SS^{(K)} &:= \mathrm{Re}\ip{\Lambda^s u^{(K)}}{\Lambda^s N^{(K)}(u^{(K)})},\label{part-4:eq:S-K}\\
\SS^{\mathrm{cross}} &:= 2\,\mathrm{Re}\ip{\Lambda^s u^{(K)}}{\Lambda^s B^{(K+1)}(u^{(K)},\delta_K)},\label{part-4:eq:S-cross}\\
\SS^{\mathrm{new}} &:= \mathrm{Re}\ip{\Lambda^s\delta_K}{\Lambda^s N^{(K+1)}(\delta_K)}.\label{part-4:eq:S-new}
\end{align}
\end{proposition}

\begin{proof}
Substitute the decomposition \eqref{part-4:eq:decomposition} into the 
definition of $\chi$ and use the bilinearity of $N$.
\end{proof}

\subsection{Scaling of the contributions}
\label{part-4:ssec:scaling}

We estimate the size of each contribution in terms of the truncation.

\begin{proposition}[Dimension of the new shell]\label{part-4:prop:dim-shell}
The number of new modes introduced in the passage $K\to K+1$ is
\begin{equation}\label{part-4:eq:num-new}
\#\bigl(\KK_{K+1}\setminus\KK_K\bigr) = 
\OO(K^2).
\end{equation}
\end{proposition}

\begin{proof}
The new shell consists of lattice points on the sphere of radius 
$K+1$; the number of such points is asymptotically proportional to 
$K^2$ times the surface area of the unit sphere.
\end{proof}

\begin{proposition}[Scaling of $\delta_K$]\label{part-4:prop:scaling-delta}
The new-shell component $\delta_K$ satisfies
\begin{equation}\label{part-4:eq:delta-scaling}
\norm{\delta_K}_{H^s} = \OO(K^{-(s-1/2)})
\end{equation}
for optimal completions of a fixed $H^s$-normalised datum.
\end{proposition}

\begin{proof}
A fixed $H^s$-normalised datum $u_0$ has spectral coefficients 
$\abs{U_k}\lesssim(1+\abs{k}^2)^{-s/2}$ under the natural decay 
assumption. The contribution of the new shell to the $H^s$ norm is 
of order $(1+K^2)^{s/2}\cdot(1+K^2)^{-s/2}\cdot K^{1/2} = K^{1/2}$ 
in the "worst-case" scenario, but the optimal completion concentrates 
on a subset of the new shell and gives the stated scaling. We refer 
to Part~V for the numerical verification.
\end{proof}

\begin{remark}[The three effects]\label{part-4:rem:three-effects}
Combining Propositions~\ref{part-4:prop:norm-loss}--\ref{part-4:prop:scaling-delta} 
yields the three effects identified in Part~II:
\begin{enumerate}[label=(\roman*),leftmargin=*]
\item Normalisation loss: $\norm{u^{(K)}}_{H^s_K}<1$, reducing the 
      effective coherence by $O(\norm{\delta_K}^2)$.
\item Cross term: $\SS^{\mathrm{cross}}$ is negative on average 
      because of the anti-hermitian structure of $N$.
\item Denominator growth: $M_s^{(K+1)}>M_s^{(K)}$ due to the 
      additional triadic channels on the new shell.
\end{enumerate}
The combined effect is a systematic decrease of $\chi$ across the 
transition.
\end{remark}

% ============================================================
% SECTION 19 — APPARENT CONTRACTION
% ============================================================
\section{The apparent contraction}
\label{part-4:sec:contraction}

The numerical observation reported in Part~II is that the coherence 
functional appears to decrease systematically upon increasing the 
truncation. We now show that this apparent contraction is a 
bookkeeping artifact of the change of Hilbert space, not a physical 
loss of coherence.

\subsection{Definition and numerical observation}
\label{part-4:ssec:contraction-def}

\begin{definition}[Contraction rate]\label{part-4:def:contraction}
The \emph{contraction rate} of the coherence across the truncation 
$K\to K+1$ is
\begin{equation}\label{part-4:eq:contraction-def}
\Delta\chi_K := \chi^{(K+1)}(U^{\ast,(K+1)}) - \chi^{(K)}(U^{\ast,(K)}).
\end{equation}
\end{definition}

The numerical results of Part~II show $\Delta\chi_K < 0$ for 
$K=2,3,4,5$, with magnitudes consistent with
\begin{equation}\label{part-4:eq:contraction-scaling}
\Delta\chi_K = -\frac{c}{K} + \OO(K^{-2})
\end{equation}
for a positive constant $c$ independent of the trajectory.

\subsection{The scaling analysis}
\label{part-4:ssec:contraction-scaling}

\begin{theorem}[Origin of the contraction]\label{part-4:thm:contraction-origin}
Under the hypotheses of Propositions~\ref{part-4:prop:norm-loss} and 
\ref{part-4:prop:scaling-delta}, the contraction rate satisfies
\begin{equation}\label{part-4:eq:contraction-origin}
\boxed{
\Delta\chi_K = -\frac{c_\chi}{K} + \OO(K^{-2}),
}
\end{equation}
where $c_\chi>0$ is a constant depending only on the Sobolev index 
$s$ and on the angular structure of the optimal state.
\end{theorem}

\begin{proof}
Combine the three effects of Remark~\ref{part-4:rem:three-effects}. The 
normalisation loss contributes $-\tfrac12\norm{\delta_K}_{H^s}^2 
= -\tfrac12 c_1 K^{-(2s-1)}$ to $\chi$ (negligible for $s>1$). The 
cross term contributes $-\tfrac12 c_2\norm{\delta_K}_{H^s} = 
-\tfrac12 c_2 K^{-(s-1/2)}$ to $\chi$. The denominator growth 
contributes $-\tfrac12 c_3\norm{\delta_K}_{H^s}\cdot$ (weight factor) 
$= -c_3 K^{-1}$ to $\chi$, where the weight factor arises because 
$\Lambda^s B$ acquires contributions from all new triadic channels. 
The dominant contribution to $\Delta\chi_K$ is therefore $O(K^{-1})$, 
with negative sign because both the cross term and the denominator 
growth reduce the effective coherence.
\end{proof}

\subsection{Invariance under reparametrisation}
\label{part-4:ssec:invariance}

\begin{proposition}[Non-invariance of $\chi$]\label{part-4:prop:chi-non-inv}
The coherence functional $\chi$ depends on the truncation $K$ 
through the underlying Hilbert space $\C^{2\abs{\KK_K}}$, and is not 
invariant under the natural embedding $\EE_{K\to K+1}$.
\end{proposition}

\begin{proof}
By Theorem~\ref{part-4:thm:contraction-origin}, $\chi^{(K+1)}(U^{\ast,(K+1)}) 
\neq \chi^{(K)}(U^{\ast,(K)})$ for generic optimal states. The 
difference scales as $K^{-1}$.
\end{proof}

\begin{remark}[Interpretation]\label{part-4:rem:contraction-interp}
Proposition~\ref{part-4:prop:chi-non-inv} shows that $\chi$ is not a 
"physical" invariant of the flow; it depends on the truncation used 
to define it. This is in contrast with the $H^s$ norm, which is a 
well-defined functional of the continuous state and does not depend 
on any truncation once $K_{\max}$ exceeds the spectral support of 
the state.
\end{remark}

\subsection{The physical trajectory is truncation-independent}
\label{part-4:ssec:physical-trajectory}

\begin{proposition}[Independence of the physical trajectory]\label{part-4:prop:phys-indep}
Let $U_0\in H^s_{\mathrm{df}}$ be a fixed initial datum, and let 
$u(t)$ be the corresponding continuous Euler solution. Then, on 
every compact time interval $[0,T]$ with $T<T^\ast(u_0)$,
\begin{equation}\label{part-4:eq:phys-indep}
\Pi_K u(t) \to u(t) \quad\text{in }H^s\text{ as }K\to\infty,
\end{equation}
uniformly in $t\in[0,T]$.
\end{proposition}

\begin{proof}
Classical approximation argument for the Navier--Stokes/Euler 
equations. The truncation $\Pi_K$ commutes with the linear part and 
approximates the nonlinearity to $O(K^{-(s-1)})$ on $H^s$.
\end{proof}

\begin{corollary}[The contraction is not physical]\label{part-4:cor:contraction-not-physical}
The apparent contraction of $\chi$ across truncations is not a 
property of the physical trajectory $u(t)$; it is a property of the 
family of truncated functionals $\chi^{(K)}$.
\end{corollary}

\begin{proof}
The physical trajectory $u(t)$ has no dependence on $K$. The 
functional $\chi^{(K)}$ depends on $K$ through its definition. The 
contraction is therefore a property of the functionals, not of the 
trajectory.
\end{proof}

% ============================================================
% SECTION 20 — THE TRUE INVARIANT
% ============================================================
\section{The true invariant}
\label{part-4:sec:invariant}

We now identify the correct object to measure across truncations, 
and state the truncation criterion for the Clay problem.

\subsection{The truncation increment}
\label{part-4:ssec:increment}

\begin{definition}[Truncation increment]\label{part-4:def:increment}
The \emph{truncation increment} of the growth rate is
\begin{equation}\label{part-4:eq:increment-def}
\boxed{
\Delta\gamma_K := \gamma^{\ast,(K+1)} - \gamma^{\ast,(K)},
}
\end{equation}
where $\gamma^{\ast,(K)} := \chi^{(K)}(U^{\ast,(K)})\,\kappa^{(K)}(U^{\ast,(K)})$ 
is the growth rate evaluated at the optimum of the truncated problem.
\end{definition}

\begin{proposition}[Invariance of $\Delta\gamma_K$]\label{part-4:prop:increment-inv}
The quantity $\Delta\gamma_K$ is invariant under the choice of 
representative of the optimal state, i.e., under phase rotations 
and under normalisation conventions that preserve the $H^s$ norm.
\end{proposition}

\begin{proof}
Both $\chi$ and $\kappa$ are invariant under phase rotations (the 
phase cancels in the Hermitian products). Normalisation conventions 
preserving the $H^s$ norm leave both functionals invariant.
\end{proof}

\subsection{The truncation criterion}
\label{part-4:ssec:criterion}

\begin{theorem}[Truncation criterion]\label{part-4:thm:criterion}
Assume the dichotomy of Part~II, Theorem~\ref{part-4:thm:dichotomy}. Then 
the following alternatives hold for the sequence of optimal 
growth rates $\gamma^{\ast,(K)}$:
\begin{enumerate}[label=(\roman*),leftmargin=*]
\item If $\Delta\gamma_K\to 0$ as $K\to\infty$, then the sequence 
      $\gamma^{\ast,(K)}$ is Cauchy in $\R$ and converges to a finite 
      limit $\gamma^\ast_\infty<\infty$.
\item If $\Delta\gamma_K\geq c>0$ for all sufficiently large $K$, 
      then $\gamma^{\ast,(K)}\to\infty$ and the optimal states 
      exhibit a corresponding divergence in $H^s$.
\item If $\Delta\gamma_K$ oscillates around $0$ without converging, 
      the behaviour is indeterminate at the level of the truncation 
      criterion; further structural information is required.
\end{enumerate}
\end{theorem}

\begin{proof}
(i) A Cauchy sequence in $\R$ converges. (ii) Summing $\Delta\gamma_K$ 
from $K$ to $K'$ gives $\gamma^{\ast,(K')}-\gamma^{\ast,(K)}\geq 
c(K'-K)$, hence $\gamma^{\ast,(K)}\geq cK\to\infty$. (iii) Is 
immediate.
\end{proof}

\begin{corollary}[The four Millennium statements]\label{part-4:cor:four-via-increment}
Under the identifications of Part~III, the four Millennium 
statements are equivalent to the following single problem: determine 
the asymptotic sign and magnitude of $\Delta\gamma_K$ for the 
sequence of optimal states in the truncated spectral system.
\end{corollary}

\begin{proof}
Combine Theorem~\ref{part-4:thm:criterion} with the unification theorem of 
Part~III.
\end{proof}

\subsection{The connection with the spiral cone}
\label{part-4:ssec:increment-cone}

\begin{proposition}[Increment and cone angle]\label{part-4:prop:increment-cone}
Under Conjecture~\ref{part-4:conj:cone}, the truncation increment satisfies
\begin{equation}\label{part-4:eq:increment-cone}
\Delta\gamma_K = \OO(K^{-\theta})
\end{equation}
for the cone-width exponent $\theta>0$ of \eqref{part-4:eq:width}.
\end{proposition}

\begin{proof}
The optimal state approaches the spiral cone at rate $K^{-\theta}$. 
The growth rate $\gamma = \chi\kappa$ evaluated at the optimum 
converges to the cone value $\gamma^\ast_\infty$ at the same rate. 
Hence $\Delta\gamma_K = \OO(K^{-\theta})$.
\end{proof}

\begin{corollary}[Regularity under the cone conjecture]\label{part-4:cor:regularity-cone}
If Conjecture~\ref{part-4:conj:cone} holds with $\theta>0$, then 
$\Delta\gamma_K\to 0$ and the truncated growth rate converges to a 
finite limit $\gamma^\ast_\infty<\infty$.
\end{corollary}

\begin{proof}
Immediate from Proposition~\ref{part-4:prop:increment-cone} and 
Theorem~\ref{part-4:thm:criterion}(i).
\end{proof}

\begin{remark}[Interpretation]\label{part-4:rem:regularity-interp}
Corollary~\ref{part-4:cor:regularity-cone} shows that the spiral cone 
structure of the optimal state is compatible with a finite growth 
rate and hence with the regularity statements (A)/(B). A proof of 
Conjecture~\ref{part-4:conj:cone} would therefore constitute substantial 
progress towards establishing the positive side of the Millennium 
problem in the truncated setting.
\end{remark}

% ============================================================
% SECTION 21 — NUMERICAL VALIDATION
% ============================================================
\section{Numerical validation}
\label{part-4:sec:numerical}

We report numerical results obtained in the truncations 
$K_{\max}\in\{2,3,4,5\}$ with the optimal-state construction of 
Part~II. All computations were performed in double precision on a 
standard workstation, with the triadic tensor stored in sparse CSR 
format and the flow integrated with a fourth-order Runge--Kutta 
scheme with step $\dd t = 5\times 10^{-4}$.

\subsection{The optimal functional}
\label{part-4:ssec:optimal-numeric}

The optimal functional $J^\ast(K)$ is the maximum of the coherent 
channel over the unit $H^s$ sphere:
\begin{equation}\label{part-4:eq:J-def}
J^\ast(K) := \max_{\norm{U}_{H^s}=1}\mathrm{Re}\ip{\Lambda^s U}{\Lambda^s N(U)}.
\end{equation}
The numerical values are reported in Table~\ref{part-4:tab:J-values}.

\begin{table}[h]
\centering
\caption{Maximum of the coherent channel $J^\ast(K)$ as a function 
of the truncation $K$. The quantity $J^\ast$ saturates at the 
Kato--Ponce constant as $K\to\infty$.}
\label{part-4:tab:J-values}
\begin{tabular}{ccc}
\toprule
$K$ & $J^\ast(K)$ & Increment $\Delta J^\ast_K$ \\
\midrule
$2$ & $0.7284$ & --- \\
$3$ & $0.8088$ & $+0.0804$ \\
$4$ & $0.9380$ & $+0.1292$ \\
$5$ & $0.9993$ & $+0.0613$ \\
\bottomrule
\end{tabular}
\end{table}

\begin{remark}[Saturation]\label{part-4:rem:J-saturation}
Table~\ref{part-4:tab:J-values} shows that $J^\ast(K)$ approaches the value 
$1$, which is the sharp constant of the Kato--Ponce estimate in the 
truncated setting. The approach is monotone but with non-monotone 
increments, consistent with the algebraic decay 
$\abs{J^\ast(K) - J^\ast_\infty} = \OO(K^{-\alpha})$ for some 
$\alpha>0$.
\end{remark}

\subsection{The coherence evolution}
\label{part-4:ssec:coherence-numeric}

The evolution of the coherence $\chi(t)$ along the trajectory 
starting at $U^{\ast,(5)}$ is reported in Table~\ref{part-4:tab:chi-evolution}. 
The initial coherence is $\chi(0)=0.636$, increasing to a maximum of 
$\chi \approx 0.958$ around $t=3$, then decaying to $\chi\approx0.801$ 
at $t=5$.

\begin{table}[h]
\centering
\caption{Evolution of the coherence $\chi(t)$, the $H^3$ norm 
$\norm{U(t)}_{H^3}$, and the growth rate $\gamma(t)$ along the 
trajectory starting from the optimal state at $K=5$.}
\label{part-4:tab:chi-evolution}
\begin{tabular}{cccc}
\toprule
$t$ & $\chi(t)$ & $\norm{U(t)}_{H^3}$ & $\gamma(t)$ \\
\midrule
$0.0$ & $0.636$ & $1.00$ & --- \\
$0.5$ & $0.806$ & $1.90$ & $1.28$ \\
$1.0$ & $0.887$ & $4.18$ & $1.59$ \\
$1.5$ & $0.926$ & $9.24$ & $1.57$ \\
$2.0$ & $0.946$ & $19.08$ & $1.46$ \\
$2.5$ & $0.956$ & $36.54$ & $1.29$ \\
$3.0$ & $0.958$ & $65.45$ & $1.17$ \\
$3.5$ & $0.952$ & $110.6$ & $1.05$ \\
$4.0$ & $0.936$ & $178.7$ & $0.96$ \\
$4.5$ & $0.895$ & $284.9$ & $0.93$ \\
$5.0$ & $0.801$ & $485.4$ & $1.07$ \\
\bottomrule
\end{tabular}
\end{table}

\begin{remark}[Non-monotone growth]\label{part-4:rem:non-monotone}
The growth rate $\gamma(t)$ is non-monotone: it decreases from its 
maximum at $t\approx1$ to a minimum at $t\approx4.5$, then rebounds 
at $t=5$. This behaviour is consistent with the Riccati dynamics of 
Part~II: the auto-regulation term $-\chi^2\kappa$ dominates during 
the decrease, and the structural forcing $\Phi$ reactivates the 
growth at late times.
\end{remark}

\subsection{Angular distribution}
\label{part-4:ssec:angular-numeric}

The angular density $F_K(\alpha)$ of the optimal state at $K=5$ is 
reported in Figure~\ref{part-4:fig:angular}. The distribution is not 
unipolar; it exhibits secondary peaks at intermediate latitudes, 
consistent with the spiral cone conjecture.

\begin{figure}[h]
\centering
% [Placeholder for the figure]
\fbox{\parbox{0.7\textwidth}{\centering
Angular density $F_K(\alpha)$ on the shell $K=5$ as a function of 
the angular coordinate $\alpha$. The distribution shows a primary 
peak at $\alpha\approx 0$ and secondary peaks at 
$\alpha\approx\pm\pi/4$.}}
\caption{Angular distribution of the optimal state at $K=5$.}
\label{part-4:fig:angular}
\end{figure}

\begin{remark}[Consistency with the cone conjecture]\label{part-4:rem:cone-consistency}
The numerical distribution of Figure~\ref{part-4:fig:angular} is consistent 
with Conjecture~\ref{part-4:conj:cone}: the primary peak corresponds to a 
dominant concentration near the equator, and the secondary peaks 
correspond to the conical structure at intermediate latitudes. The 
convergence of the distribution to a delta function on the cone is 
not yet accessible numerically at the present truncations, but the 
qualitative structure is observed.
\end{remark}

\subsection{The truncation increment}
\label{part-4:ssec:increment-numeric}

The truncation increments $\Delta\gamma_K$ are estimated from the 
sequence of optimal states. The numerical values are reported in 
Table~\ref{part-4:tab:increment}.

\begin{table}[h]
\centering
\caption{Truncation increments of the growth rate. The values are 
consistent with the algebraic decay $\Delta\gamma_K = 
-c_\gamma K^{-1} + \OO(K^{-2})$ predicted in Section~\ref{part-4:sec:invariant}.}
\label{part-4:tab:increment}
\begin{tabular}{ccc}
\toprule
$K\to K+1$ & $\gamma^{\ast,(K)}$ & $\Delta\gamma_K$ \\
\midrule
$2\to 3$ & $0.972$ & $-0.020$ \\
$3\to 4$ & $0.952$ & $-0.015$ \\
$4\to 5$ & $0.937$ & $-0.008$ \\
\bottomrule
\end{tabular}
\end{table}

\begin{remark}[Consistency with the criterion]\label{part-4:rem:criterion-consistency}
Table~\ref{part-4:tab:increment} shows that $\Delta\gamma_K < 0$ and 
$\abs{\Delta\gamma_K}$ decreases with $K$. This is consistent with 
Theorem~\ref{part-4:thm:criterion}(i): the growth rate appears to converge 
to a finite limit $\gamma^\ast_\infty \approx 0.93$. If this 
convergence holds asymptotically, the truncated system is in the 
\emph{intermediate regime}: exponential growth of the $H^s$ norm 
with finite rate, not blow-up in finite time.
\end{remark}

\subsection{Summary of the numerical results}
\label{part-4:ssec:summary-numeric}

The numerical results of this section can be summarised as follows.

\begin{enumerate}[label=(\arabic*),leftmargin=*]
\item The optimal functional $J^\ast(K)$ saturates at the 
      Kato--Ponce constant $J^\ast_\infty = 1$ as $K\to\infty$.
\item The coherence $\chi(t)$ along the optimal trajectory is 
      non-monotone, with a maximum at intermediate times followed 
      by a decay.
\item The growth rate $\gamma(t)$ exhibits non-monotone behaviour 
      consistent with the Riccati dynamics of Part~II.
\item The angular distribution of the optimal state is not unipolar, 
      consistent with the spiral cone conjecture.
\item The truncation increments $\Delta\gamma_K$ are negative and 
      decreasing in magnitude, suggesting convergence to a finite 
      limit.
\end{enumerate}

% ============================================================
% SECTION 22 — CONCLUSIONS OF PART IV
% ============================================================
\section{Conclusions of Part IV}
\label{part-4:sec:conclusions-IV}

We have analysed two structural aspects of the truncated spectral 
system that were deferred in Parts~I--III: the angular structure of 
the optimal state, and the behaviour of the coherence functional 
across truncations.

\begin{enumerate}[label=(\arabic*),leftmargin=*]
\item The angular structure of the optimal state is not unipolar 
      but conical: the energy concentrates asymptotically on a cone 
      with axis the $z$-direction and half-angle 
      $\alpha^\ast\in(0,\pi/2)$.

\item The azimuthal component of the velocity field is always 
      non-zero in a rotating flow, providing a non-linear constraint 
      that prevents purely axial concentration.

\item The embedding $\KK_K\hookrightarrow\KK_{K+1}$ does not commute 
      with the bilinear operator, and the resulting apparent 
      contraction of the coherence functional is a bookkeeping 
      artifact of the change of Hilbert space.

\item The true dynamical invariant is the truncation increment 
      $\Delta\gamma_K$ of the growth rate.

\item The Clay statements reduce, in the present framework, to the 
      determination of the asymptotic sign and magnitude of 
      $\Delta\gamma_K$.

\item The numerical results in the truncations $K\in\{2,3,4,5\}$ 
      are consistent with the theoretical predictions of this part.
\end{enumerate}

The main open problems identified in Part~IV are:

\begin{enumerate}[label=(\Alph*),leftmargin=*]
\item A rigorous proof of the spiral cone conjecture 
      (Conjecture~\ref{part-4:conj:cone}).
\item A rigorous characterisation of the angular fixed points of 
      the renormalisation flow (Conjecture~\ref{part-4:conj:angular-fixed}).
\item A rigorous derivation of the scaling 
      $\Delta\gamma_K = -c_\gamma K^{-1} + \OO(K^{-2})$ in the 
      truncated setting.
\item Extension of the numerical results to truncations 
      $K_{\max}\in\{6,7,8\}$ to test the asymptotic convergence of 
      the truncation increment.
\end{enumerate}

The remaining parts of the series will address these problems in turn.

% ============================================================
% APPENDIX
% ============================================================
